\chapter{State, Inputs, and Timing}

\section{Nominal State}

The nominal state is nonlinear:
\begin{equation}
\xb_n =
\left(
\pb_{eb}^{e},
\vb_{eb}^{e},
\quat_b^e,
\bb_g,
\bb_a
\right),
\end{equation}
where $\bb_g$ and $\bb_a$ are the estimated gyro and accelerometer biases. The
quaternion is not stored in the linear Kalman error state as four independent
components, but rather is represented in the filter as a three-component
small-angle perturbation about the nominal quaternion attitude state.

\section{Error State}

The first error-state vector is
\begin{equation}
\deltax =
\begin{bmatrix}
\delta\pb^{e}\\
\delta\vb^{e}\\
\delta\boldsymbol{\theta}\\
\delta\bb_g\\
\delta\bb_a
\end{bmatrix}
\in \R^{15}.
\end{equation}

The attitude error $\delta\boldsymbol{\theta}$ is a small-angle perturbation
resolved in ECEF and defined by \cref{eq:v1_attitude_error_dcm_relation}. The
same convention must be used by propagation, measurement Jacobians, injection,
and reset.

\section{Rotation-Vector Maps}

The implementation will repeatedly map small rotation vectors into attitude
updates. For a rotation vector $\boldsymbol\alpha\in\R^3$, define
$\alpha=\norm{\boldsymbol\alpha}$ and, for $\alpha>0$,
$\mathbf{u}=\boldsymbol\alpha/\alpha$. Rodrigues' formula gives the direction
cosine matrix
\begin{equation}
C(\boldsymbol\alpha)
=
\I
+
\frac{\sin\alpha}{\alpha}\crossmat{\boldsymbol\alpha}
+
\frac{1-\cos\alpha}{\alpha^2}
\crossmat{\boldsymbol\alpha}^2,
\label{eq:v1_rodrigues_dcm}
\end{equation}
with the small-angle limit
\begin{equation}
C(\boldsymbol\alpha)
\approx
\I+\crossmat{\boldsymbol\alpha}.
\end{equation}
The corresponding scalar-first quaternion is
\begin{equation}
q(\boldsymbol\alpha)
=
\begin{bmatrix}
\cos\left(\alpha/2\right)\\
\mathbf{u}\sin\left(\alpha/2\right)
\end{bmatrix},
\label{eq:v1_rotation_vector_quaternion}
\end{equation}
with small-angle limit
\begin{equation}
q(\boldsymbol\alpha)
\approx
\begin{bmatrix}
1\\
\frac{1}{2}\boldsymbol\alpha
\end{bmatrix}.
\end{equation}
The sign of $\boldsymbol\alpha$ must follow the same attitude-error convention
as \cref{eq:v1_attitude_error_dcm_relation}; the maps above define the
rotation-vector-to-attitude primitives used by mechanization, correction, and
tests \cite{groves2013,sola2018}.

\section{Continuous-Time IMU Model}

Before discretization into increments, the basic v1 IMU measurement model is
\begin{align}
\tilde{\omegab}_{ib}^{b}
&=
\omegab_{ib}^{b}
+
\bb_g
+
\mathbf{w}_g,\\
\tilde{\fb}_{ib}^{b}
&=
\fb_{ib}^{b}
+
\bb_a
+
\mathbf{w}_a,\\
\dot{\bb}_g
&=
\mathbf{w}_{bg},\\
\dot{\bb}_a
&=
\mathbf{w}_{ba}.
\end{align}
Scale factors, misalignment, $g$-sensitivity, and other higher-order IMU
calibration terms are intentionally deferred, but this section is the expected
home for those terms when the model grows.

\section{IMU Input Contract}
\label{sec:imu_input_contract}

The first implementation consumes timestamped IMU increments with:
\begin{align}
\Delta\tilde{\boldsymbol\theta}_{ib}^{b} & && \text{measured body wrt inertial angle increment, resolved in body},\\
\Delta\tilde{\mathbf{v}}_{ib}^{b} & && \text{measured specific-force velocity increment, resolved in body},\\
t_k & && \text{sample time}.
\end{align}

The simulator may generate these increments by integrating ideal angular rate
and specific force over the sample interval, but the mechanization contract is
increment-based. Bias correction is applied before coning/sculling compensation:
\begin{align}
\Delta\hat{\boldsymbol\theta}_{ib}^{b}
&=
\Delta\tilde{\boldsymbol\theta}_{ib}^{b}
-
\hat{\bb}_g\Delta t,\\
\Delta\hat{\mathbf{v}}_{ib}^{b}
&=
\Delta\tilde{\mathbf{v}}_{ib}^{b}
-
\hat{\bb}_a\Delta t.
\end{align}
The simple $\hat{\bb}\Delta t$ correction assumes the estimated bias is constant
over the interval, or equivalently that $\hat{\bb}$ represents the interval
average used by the discrete mechanization. Scale factors and misalignments are
deferred from v1, but the interface should not make them impossible to add.

\section{Timing Architecture}

The Navigator API should support three time scales:
\begin{itemize}
    \item high-rate strapdown integration from IMU data, typically
    \SIrange{600}{1000}{\hertz};
    \item medium-rate covariance/state-transition propagation, typically
    \SIrange{100}{200}{\hertz} or event/time driven;
    \item low-rate asynchronous aiding updates, such as GNSS, barometer, or
    future sensors.
\end{itemize}

The implementation should buffer enough timing information to align aiding
measurements to the propagated state. The v1 implementation should include
fixed-capacity high-rate IMU increment buffers and medium-rate nominal-state
and STM history buffers. Those buffers are stored for validation, profiling, and
future latency support; v1 does not yet extrapolate delayed aiding measurements
to the current time. Delayed-measurement extrapolation is a v2 feature.

The first version may keep the buffering simple, but ownership must be explicit:
IMU samples feed the mechanization path, state-transition data feeds the
filter-prediction/history path, and aiding measurements feed the update path.
